\documentclass[11pt]{article}
\usepackage[margin=1in]{geometry}
\usepackage{amsmath,amssymb,amsthm,booktabs,longtable,array,enumitem}

\newtheorem{definition}{Definition}
\newtheorem{theorem}{Theorem}
\newtheorem{proposition}{Proposition}
\newtheorem{corollary}{Corollary}
\newtheorem{obstruction}{Precise obstruction}
\newenvironment{nonconclusion}{\par\smallskip\noindent\textbf{Non-conclusion.}\ }{\par\smallskip}

\newcommand{\Src}{\mathrm{src}}
\newcommand{\Bool}{\mathbb B}
\newcommand{\Rel}{\mathcal R}
\newcommand{\Reach}{\mathcal Q}
\newcommand{\Clock}{\mathsf{Clock}_{\Src}}
\newcommand{\Rod}{\mathsf{Rod}_{\Src}}
\newcommand{\Signal}{\mathsf{Signal}_{\Src}}
\newcommand{\Detector}{\mathsf{Detector}_{\Src}}
\newcommand{\FreeFall}{\mathsf{FreeFall}_{\Src}}

\title{CommonSourcePropagationProfileCandidate\_v1\\
\large A proposal-only finite source relation family and scoped universality obstruction}
\author{The Aether-Flow Research Control Project}
\date{2026-07-28}

\begin{document}
\maketitle

\section{Status and exact question}

This task-local object is \texttt{draft/control}, \texttt{proposal-only},
and source-extension data.  It executes v21 P7-T04 under
\texttt{CB-V21-P7-T04-COMMON-SOURCE-PROPAGATION-001}.  It does not modify
canonical ontology and it does not adopt a source propagation law.

P7-T02 supplies finite controls and nonnegative rational bookkeeping
kernels.  P7-T03 supplies five guarded protocol roles with explicit finite
source transitions and readouts.  Neither predecessor supplies physical
time, a differential equation, a covector domain, a principal symbol, a
characteristic variety, a causal cone, a metric, detector semantics, or a
matter coupling.  The exact P7-T04 question is therefore:

\begin{quote}
Which finite support and reachability relations follow from the declared
source transitions, which protocols are isomorphic as such relations under
explicit source-address bijections, and what exact invariant obstructs one
bijectively common relation for the present five-protocol suite?
\end{quote}

The answer has two parts.  A common three-address chain is constructible for
the \(\Rod\), \(\Signal\), and \(\Detector\) subfamily after one explicit
rod--signal address bijection is stipulated.  The full five-role suite does
not admit one relation up to bijective transport: already the two-address
\(\Clock\) and baseline \(\FreeFall\) controls have incompatible strongly
connected-component and directed-cycle invariants.

\section{Fixed source basis}

Let a P7-T02 control be
\[
  C=(\Omega,E,s,t,K_C),
\]
where every finite row \(E_x\) contains its typed identity
\(\mathbf 1_x\).  The coefficients \(K_C(T\mid x)\) are bookkeeping
weights.  The P7-T02 composition index \(n\in\mathbb N_0\) is a formal
update count only.

The P7-T03 protocol suite is held fixed at proposal-only status.  Its role
names \(\Clock,\Rod,\Signal,\Detector,\FreeFall\) name formal source
protocols only.  In particular, ``signal'', ``detector'', and ``free fall''
do not import their familiar physical meanings.

\section{Finite source support and reachability}

\begin{definition}[Full and reduced support]
For a finite control \(C\), define its row-source Boolean support relation
\(\Rel_C\subseteq\Omega\times\Omega\) by
\[
  x\,\Rel_C\,y
  \quad\Longleftrightarrow\quad
  \exists T\in E_x\ \text{with}\ t(T)=y.
\]
Its adjacency matrix is
\[
  A_C(x,y)=
  \begin{cases}
  1,&x\,\Rel_C\,y,\\
  0,&\text{otherwise}.
  \end{cases}
\]
The reduced relation \(\Rel_C^\circ\) deletes only the declared typed
identity transitions.  It does not delete a nonidentity transition merely
because its source and target labels happen to agree.
\end{definition}

\begin{definition}[Boolean reachability]
For \(n\geq 0\), let
\[
  \Reach_C^{[0]}=I_\Omega,\qquad
  \Reach_C^{[n+1]}=\Reach_C^{[n]}\odot A_C,
\]
where \(\odot\) is Boolean matrix multiplication.  Thus
\(\Reach_C^{[n]}(x,y)=1\) precisely when a declared length-\(n\) transition
word connects \(x\) to \(y\).  Let
\(\Reach_C^{[\le n]}=\bigvee_{j=0}^{n}\Reach_C^{[j]}\).
\end{definition}

\begin{nonconclusion}
The integer \(n\) is not physical time, duration, affine parameter, or
energy scale.  The relation records finite source syntax and nothing more.
\end{nonconclusion}

\begin{definition}[Support profile]
The finite support profile is
\[
  \Pi(C)=
  \bigl(
  |\Omega|,\ |\Rel_C^\circ|,\ \operatorname{scc}(A_C),\
  \operatorname{cyc}(\Rel_C^\circ),\
  \operatorname{diam}_{\mathrm{reach}}(A_C)
  \bigr),
\]
where \(\operatorname{scc}\) is the number of strongly connected
components, \(\operatorname{cyc}\) records whether the reduced relation has
a directed cycle, and the reachability diameter is the largest finite
shortest directed-path length between reachable distinct vertices.
\end{definition}

\section{Transport theorem}

\begin{definition}[Compatible relation transport]
Controls \(C\) and \(C'\) are support-transport equivalent when there is a
bijection \(b:\Omega\to\Omega'\) such that
\[
  x\,\Rel_C\,y
  \quad\Longleftrightarrow\quad
  b(x)\,\Rel_{C'}\,b(y),
\]
and \(b\) carries typed identities to typed identities.  If \(P_b\) is the
permutation matrix of \(b\), this is
\[
  A_{C'}=P_b^{-1}A_CP_b
\]
over the Boolean semiring, with the convention fixed by the row-source
matrices above.
\end{definition}

\begin{theorem}[Boolean support transport]
\label{thm:transport}
If \(b\) support-transports \(C\) to \(C'\), then for every \(n\geq0\)
\[
  \Reach_{C'}^{[n]}
  =P_b^{-1}\Reach_C^{[n]}P_b,
  \qquad
  \Reach_{C'}^{[\le n]}
  =P_b^{-1}\Reach_C^{[\le n]}P_b.
\]
Consequently \(\Pi(C)=\Pi(C')\).
\end{theorem}

\begin{proof}
The first equality at \(n=1\) is the defining intertwining relation.  The
identity matrix is preserved at \(n=0\).  Induction and associativity of
Boolean matrix multiplication give the exact-length result.  Boolean union
gives the bounded-reachability result.  Vertex count, edge count, strong
connectivity, directed cycles, and shortest directed-path lengths are
invariants of relation isomorphism.
\end{proof}

\section{Explicit P7-T03 support matrices}

All matrices below use the displayed address order and row-source
convention.

\subsection{Three-address chain}

The rod, signal, and detector controls each have an identity-plus-forward
chain, hence the matrix
\[
 A_{\mathrm{ch},3}=
 \begin{pmatrix}
 1&1&0\\
 0&1&1\\
 0&0&1
 \end{pmatrix}.
\]
The reduced relation is \(0\to1\to2\), and
\[
 \Pi(C_{\mathrm{ch},3})=(3,2,3,\mathrm{false},2).
\]
P7-T03 already declares the signal--detector bijection
\(b_{SD}(s_j)=q_j\).  This P7-T04 candidate additionally stipulates the
proposal-only source-address bijection \(b_{RS}(r_j)=s_j\).  No physical
distance, timing, or detector calibration enters either bijection.

\begin{theorem}[Common three-address source skeleton]
\label{thm:subfamily}
Under \(b_{RS}\) and the P7-T03 bijection \(b_{SD}\), the \(\Rod\),
\(\Signal\), and \(\Detector\) support and bounded-reachability relations are
all transported copies of \(A_{\mathrm{ch},3}\).  In particular, address
\(0\) first reaches address \(2\) after two nonidentity forward hops.
\end{theorem}

\begin{proof}
Each declared row has exactly its typed identity and, at addresses zero and
one, the displayed forward transition.  Address two has only its identity.
The two bijections carry identities and forward transitions to the
corresponding identities and forward transitions.  Theorem~\ref{thm:transport}
then applies.
\end{proof}

\begin{nonconclusion}
The least formal hop count two is not a travel time or distance.  The common
chain is not a causal cone, metric, light propagation law, detector
universality theorem, or matter-coupling result.
\end{nonconclusion}

\subsection{Two-address clock and baseline free-fall controls}

For address order \((c_0,c_1)\), the P7-T03 clock rows give
\[
 A_{\mathrm{clk}}=
 \begin{pmatrix}
 1&1\\
 1&1
 \end{pmatrix},
 \qquad
 \Pi(C_{\mathrm{clk}})=(2,2,1,\mathrm{true},1).
\]
For address order \((f_0,f_1)\), the baseline free-fall-labelled control
gives
\[
 A_{\mathrm{ff},0}=
 \begin{pmatrix}
 1&1\\
 0&1
 \end{pmatrix},
 \qquad
 \Pi(C_{\mathrm{ff},0})=(2,1,2,\mathrm{false},1).
\]
The blocked intervention gives
\[
 A_{\mathrm{ff},\mathrm{block}}=
 \begin{pmatrix}
 1&0\\
 0&1
 \end{pmatrix}.
\]
These are source relations only.  The clock cycle is not periodic physical
motion, and the baseline relation is not geodesic motion.

\section{Minimal multi-relation countermodel}

\begin{obstruction}[Full-suite bijective universality is not derived]
\label{obst:universal}
Let the full-suite universality claim \(U_{\mathrm{bij}}\) assert that one
finite relation \((X,\Rel_\star)\) and bijections
\(b_D:\Omega_D\to X\) intertwine the support relation of every P7-T03 role
\(D\in\{\Clock,\Rod,\Signal,\Detector,\FreeFall\}\).
Then the current proposal-only suite falsifies \(U_{\mathrm{bij}}\).
\end{obstruction}

\begin{proof}
It is enough to compare the two-vertex clock and baseline free-fall
controls, so cardinality is held fixed.  The clock relation is strongly
connected and its reduced relation contains the directed two-cycle
\(c_0\to c_1\to c_0\).  The baseline free-fall relation has two strongly
connected components and its reduced relation is the acyclic single edge
\(f_0\to f_1\).  Strongly connected-component count and directed-cycle
content are invariant under every bijection.  Therefore no bijection
intertwines these two relations.  A fortiori, no one relation can be
bijectively common to all five protocols.
\end{proof}

The obstruction identifier is
\[
 \texttt{OBST-P7T04-FULL-SUITE-COMMON-PROPAGATION-001}.
\]
Its scope is exact: it rejects \(U_{\mathrm{bij}}\) for the present explicit
P7-T03 controls.  It does not reject every possible weaker homomorphism,
coarse quotient, enriched relation, stochastic comparison, continuum limit,
or future source extension.  A one-point quotient would be formally common
but would erase every discriminating relation; the current records neither
select it nor give it physical meaning.

\section{Why no principal symbol, cone, or metric follows}

A principal-symbol calculation requires at least a declared differential or
pseudodifferential operator, a leading-order term, and a covector variable
\(\xi\) over a base space.  A characteristic set requires a polynomial or
equivalent symbol condition in that covector.  A cone or metric
factorization requires still more structure: dimension, differentiable
locality, scale behavior, signature or causal interpretation, and a map from
source records to those objects.

The present finite relations supply none of:
\begin{enumerate}[leftmargin=*]
\item a source manifold, tangent or cotangent bundle, or covector \(\xi\);
\item a local differential operator or leading-order derivative;
\item a characteristic polynomial homogeneous in \(\xi\);
\item an energy, wavelength, momentum, or low-energy scale variable;
\item a map from Boolean reachability to a physical causal or metric
  structure; or
\item a derivation tying all matter sectors to one action or coupling.
\end{enumerate}

One can form an ordinary matrix characteristic polynomial
\(\det(\lambda I-A_C)\), but \(\lambda\) is then only a spectral parameter of
a finite adjacency matrix.  Calling it a dispersion relation or identifying
its zeros with a physical characteristic variety would add an unauthorized
map.  Thus P7-T04 returns a precise source-level obstruction before any
principal-symbol or common-cone claim is well typed.

\begin{nonconclusion}
The absence of those maps from the current candidate does not prove that a
future conservative source extension cannot supply them.  It proves only
that they are not derived here.
\end{nonconclusion}

\section{Universality conditions and reopening criteria}

A future packet may reopen the same physical burden only with materially new
source-side data satisfying all of the following:
\begin{enumerate}[leftmargin=*]
\item one declared propagation carrier or relation family for every admitted
  matter protocol;
\item explicit comparison maps with a nontrivial universality criterion;
\item stability of that criterion under allowed source presentations and
  interventions;
\item an independently justified scale or limiting parameter if
  ``low-energy'' is asserted;
\item a source-to-operator map sufficient to define a principal symbol, or a
  precise source analogue with its own interpretation theorem; and
\item a separate model-to-world map before any physical cone, metric,
  detector, or coupling language is licensed.
\end{enumerate}

Merely renaming Boolean reachability, dropping incompatible protocols,
quotienting every relation to one point, or importing the desired metric into
each matter action does not reopen the obstruction.

\section{Assumption and authority ledger}

The new proposal inputs are the support-profile definition and the
rod--signal address bijection \(b_{RS}\).  The transport theorem, the common
three-address subfamily theorem, and the two-vertex countermodel follow
exactly from those inputs and the fixed P7-T03 rows.

The current ontology does not derive or adopt the support-profile choice,
\(b_{RS}\), a universal relation, a principal-symbol map, or a physical
interpretation.  Adoption remains
\texttt{blocked\_adoption\_open\_continuation}.  Validator success, the
AgentJob, the checkpoint, and this document's recency are process evidence,
not proof or promotion authority.

Forbidden conclusions include canonical ontology or source-law adoption,
physical time or causality, signal speed, a common cone or metric, empirical
detector meaning, universal coupling, equivalence-principle behavior, matter
action, stress energy, Einstein equations, exact-GR benchmark recovery,
publication, completed derivation, global theory rejection, and future
source-extension impossibility.

\end{document}
